\section{Appendix}

The complete source code for Lab 6.1 is organised below by architectural layer. Every module is reproduced here in full so that the report can be read as a self-contained reference; the latest version always lives in the project repository.

\textbf{GitHub Repository:} \url{https://github.com/mcittkmims/es-labs}

\subsection{Application Layer}

\subsubsection{Lab Selector --- src/main.cpp}

The central entry point delegates \texttt{setup()} and \texttt{loop()} to the active lab through preprocessor flags set in \texttt{platformio.ini}; only the Lab 6.1 branch is shown here for brevity.

\begin{lstlisting}[language=C++, caption=src/main.cpp --- Lab selector (excerpt for Lab 6.1), label=lst:appendix_main]
#include <Arduino.h>

#if defined(LAB6_1)
    #include "lab6_1_main.h"
#endif

void setup() {
#if defined(LAB6_1)
    lab6_1Setup();
#endif
}

void loop() {
#if defined(LAB6_1)
    lab6_1Loop();
#endif
}
\end{lstlisting}

\subsubsection{Lab 6.1 Entry Point}

\begin{lstlisting}[language=C++, caption=lab/lab6\_1/lab6\_1\_main.h --- public interface, label=lst:appendix_lab6_1_h]
#ifndef LAB6_1_MAIN_H
#define LAB6_1_MAIN_H

void lab6_1Setup();
void lab6_1Loop();

#endif // LAB6_1_MAIN_H
\end{lstlisting}

\begin{lstlisting}[language=C++, caption=lab/lab6\_1/lab6\_1\_main.cpp --- application orchestration, label=lst:appendix_lab6_1_cpp]
#include "lab6_1_main.h"
#include <Arduino.h>
#include <stdio.h>

#include "Button.h"
#include "ButtonLedFsm.h"
#include "Led.h"
#include "StdioSerial.h"

static const uint8_t       PIN_BUTTON         = 7;
static const uint8_t       PIN_LED            = 8;
static const unsigned long BAUD_RATE          = 9600;
static const uint16_t      BUTTON_DEBOUNCE_MS = 30;
static const uint16_t      LOOP_TICK_MS       = 5;

static Button       button(PIN_BUTTON, /*activeLow=*/true, BUTTON_DEBOUNCE_MS);
static Led          led(PIN_LED);
static ButtonLedFsm fsm;

static void printFsmTable() {
    printf("FSM state table (Moore, 2 states):\r\n");
    printf("  Num | Name    | Out | In=0    | In=1    \r\n");
    printf("  ----+---------+-----+---------+---------\r\n");
    printf("  0   | LED_OFF |  0  | LED_OFF | LED_ON  \r\n");
    printf("  1   | LED_ON  |  1  | LED_ON  | LED_OFF \r\n");
}

void lab6_1Setup() {
    stdioSerialInit(BAUD_RATE);
    button.init();
    led.init();
    fsm.init();
    led.set(fsm.getOutput() != 0);

    printf("\r\n========================================\r\n");
    printf("  Lab 6.1: Button-LED Finite State        \r\n");
    printf("           Machine (Moore, 2 states)     \r\n");
    printf("  MCU: Arduino Mega 2560                  \r\n");
    printf("========================================\r\n");
    printf("Wiring:\r\n");
    printf("  Button -> D%u (active-LOW, INPUT_PULLUP)\r\n", PIN_BUTTON);
    printf("  LED    -> D%u (220 Ohm to GND)\r\n",            PIN_LED);
    printf("Debounce window: %u ms\r\n", (unsigned)BUTTON_DEBOUNCE_MS);
    printFsmTable();
    printf("[FSM] Initial state: %s\r\n", fsm.getStateName());
    printf("Press the button to toggle the LED state.\r\n\r\n");
    fsm.clearChanged();
}

void lab6_1Loop() {
    button.update();
    if (button.wasPressed()) {
        fsm.processEvent();
    }
    if (fsm.changed()) {
        led.set(fsm.getOutput() != 0);
        printf("[FSM] State -> %s (LED %s)\r\n",
               fsm.getStateName(),
               led.isOn() ? "ON" : "OFF");
        fsm.clearChanged();
    }
    delay(LOOP_TICK_MS);
}
\end{lstlisting}

\subsection{Service Layer}

\subsubsection{ButtonLedFsm Module}

\begin{lstlisting}[language=C++, caption=lib/ButtonLedFsm/ButtonLedFsm.h --- two-state Moore FSM interface, label=lst:appendix_fsm_h]
#ifndef BUTTON_LED_FSM_H
#define BUTTON_LED_FSM_H

#include <stdint.h>

enum ButtonLedState {
    LED_OFF_STATE = 0,
    LED_ON_STATE  = 1,
    BUTTON_LED_STATE_COUNT
};

class ButtonLedFsm {
public:
    ButtonLedFsm();
    void           init();
    void           processEvent();
    ButtonLedState getState() const;
    uint8_t        getOutput() const;
    const char    *getStateName() const;
    bool           changed() const;
    void           clearChanged();
private:
    struct StateRow {
        uint8_t        output;
        ButtonLedState next[2];
    };
    static const StateRow _table[BUTTON_LED_STATE_COUNT];
    ButtonLedState _state;
    bool           _changed;
};

#endif // BUTTON_LED_FSM_H
\end{lstlisting}

\begin{lstlisting}[language=C++, caption=lib/ButtonLedFsm/ButtonLedFsm.cpp --- FSM implementation, label=lst:appendix_fsm_cpp]
#include "ButtonLedFsm.h"

const ButtonLedFsm::StateRow ButtonLedFsm::_table[BUTTON_LED_STATE_COUNT] = {
    { 0, { LED_OFF_STATE, LED_ON_STATE  } },
    { 1, { LED_ON_STATE,  LED_OFF_STATE } }
};

ButtonLedFsm::ButtonLedFsm()
    : _state(LED_OFF_STATE), _changed(true) {}

void ButtonLedFsm::init() {
    _state   = LED_OFF_STATE;
    _changed = true;
}

void ButtonLedFsm::processEvent() {
    ButtonLedState next = _table[_state].next[1];
    if (next != _state) {
        _state   = next;
        _changed = true;
    }
}

ButtonLedState ButtonLedFsm::getState() const {
    return _state;
}

uint8_t ButtonLedFsm::getOutput() const {
    return _table[_state].output;
}

const char *ButtonLedFsm::getStateName() const {
    return (_state == LED_ON_STATE) ? "ON" : "OFF";
}

bool ButtonLedFsm::changed() const  { return _changed; }
void ButtonLedFsm::clearChanged()   { _changed = false; }
\end{lstlisting}

\subsection{ECAL Layer}

\subsubsection{StdioSerial Module}

\begin{lstlisting}[language=C++, caption=lib/StdioSerial/StdioSerial.h --- printf-to-UART bridge, label=lst:appendix_stdio_h]
#ifndef STDIO_SERIAL_H
#define STDIO_SERIAL_H

#include <Arduino.h>
#include <stdio.h>

void stdioSerialInit(unsigned long baudRate);

#endif // STDIO_SERIAL_H
\end{lstlisting}

\begin{lstlisting}[language=C++, caption=lib/StdioSerial/StdioSerial.cpp --- bridge implementation, label=lst:appendix_stdio_cpp]
#include "StdioSerial.h"

static int serialPutChar(char c, FILE *stream) {
    Serial.write(c);
    return 0;
}

static int serialGetChar(FILE *stream) {
    while (!Serial.available()) { }
    char c = Serial.read();
    if (c == '\r') {
        Serial.write('\r');
        Serial.write('\n');
        return '\n';
    }
    if (c == '\b' || c == 127) {
        Serial.write('\b');
        Serial.write(' ');
        Serial.write('\b');
        return c;
    }
    Serial.write(c);
    return c;
}

static FILE serialStream;

void stdioSerialInit(unsigned long baudRate) {
    Serial.begin(baudRate);
    while (!Serial) { ; }
    fdev_setup_stream(&serialStream,
        serialPutChar, serialGetChar, _FDEV_SETUP_RW);
    stdout = &serialStream;
    stdin  = &serialStream;
}
\end{lstlisting}

\subsection{MCAL Layer}

\subsubsection{Button Driver}

\begin{lstlisting}[language=C++, caption=lib/Button/Button.h --- debounced push-button driver interface, label=lst:appendix_button_h]
#ifndef BUTTON_H
#define BUTTON_H

#include <Arduino.h>

class Button {
public:
    Button(uint8_t pin, bool activeLow = true, uint16_t debounceMs = 30);
    void init();
    void update();
    bool isPressed() const;
    bool wasPressed();
    bool wasReleased();
private:
    uint8_t  _pin;
    bool     _activeLow;
    uint16_t _debounceMs;
    bool     _stableState;
    bool     _lastRawState;
    uint32_t _lastChangeMs;
    bool     _pressedFlag;
    bool     _releasedFlag;
    bool     readLogical() const;
};

#endif // BUTTON_H
\end{lstlisting}

\begin{lstlisting}[language=C++, caption=lib/Button/Button.cpp --- driver implementation, label=lst:appendix_button_cpp]
#include "Button.h"

Button::Button(uint8_t pin, bool activeLow, uint16_t debounceMs)
    : _pin(pin), _activeLow(activeLow), _debounceMs(debounceMs),
      _stableState(false), _lastRawState(false),
      _lastChangeMs(0), _pressedFlag(false), _releasedFlag(false) {}

void Button::init() {
    if (_activeLow) {
        pinMode(_pin, INPUT_PULLUP);
    } else {
        pinMode(_pin, INPUT);
    }
    _lastRawState = readLogical();
    _stableState  = _lastRawState;
    _lastChangeMs = millis();
    _pressedFlag  = false;
    _releasedFlag = false;
}

void Button::update() {
    bool raw = readLogical();
    uint32_t now = millis();
    if (raw != _lastRawState) {
        _lastRawState = raw;
        _lastChangeMs = now;
        return;
    }
    if (raw != _stableState && (now - _lastChangeMs) >= _debounceMs) {
        _stableState = raw;
        if (_stableState) _pressedFlag  = true;
        else              _releasedFlag = true;
    }
}

bool Button::isPressed() const  { return _stableState; }
bool Button::wasPressed()       {
    bool flag = _pressedFlag;  _pressedFlag = false;  return flag; }
bool Button::wasReleased()      {
    bool flag = _releasedFlag; _releasedFlag = false; return flag; }

bool Button::readLogical() const {
    int level = digitalRead(_pin);
    return _activeLow ? (level == LOW) : (level == HIGH);
}
\end{lstlisting}

\subsubsection{LED Driver}

\begin{lstlisting}[language=C++, caption=lib/Led/Led.h --- LED driver interface, label=lst:appendix_led_h]
#ifndef LED_H
#define LED_H

#include <Arduino.h>

class Led {
public:
    Led(uint8_t pin);
    void init();
    void turnOn();
    void turnOff();
    void toggle();
    void set(bool on);
    bool isOn() const;
private:
    uint8_t ledPin;
    bool    state;
};

#endif // LED_H
\end{lstlisting}

\begin{lstlisting}[language=C++, caption=lib/Led/Led.cpp --- LED driver implementation, label=lst:appendix_led_cpp]
#include "Led.h"

Led::Led(uint8_t pin) : ledPin(pin), state(false) {}

void Led::init() {
    pinMode(ledPin, OUTPUT);
    digitalWrite(ledPin, LOW);
    state = false;
}

void Led::turnOn() {
    digitalWrite(ledPin, HIGH);
    state = true;
}

void Led::turnOff() {
    digitalWrite(ledPin, LOW);
    state = false;
}

void Led::toggle() {
    if (state) turnOff(); else turnOn();
}

void Led::set(bool on) {
    if (on) turnOn(); else turnOff();
}

bool Led::isOn() const { return state; }
\end{lstlisting}

\subsection{Build and Simulation Configuration}

\begin{lstlisting}[caption=labs/platformio.ini --- Lab 6.1 environment, label=lst:appendix_pio_ini]
[env:lab6_1]
platform     = atmelavr
board        = megaatmega2560
framework    = arduino
monitor_speed = 9600
build_src_filter = +<*> +<../lab/lab6_1/*>
build_flags  = -I lab/lab6_1 -DLAB6_1
\end{lstlisting}

\begin{lstlisting}[caption=labs/wokwi/lab6.1/wokwi.toml --- simulator firmware path, label=lst:appendix_wokwi_toml]
[wokwi]
version  = 1
firmware = "../../.pio/build/lab6_1/firmware.hex"
elf      = "../../.pio/build/lab6_1/firmware.elf"
\end{lstlisting}
